\section{An infinite exact failure at five colors and its six-color repair}
\label{sec:blind}

\subsection{The weaving family as an adversarial test}

Let
\[
 W_m=\widehat{(\sigma_1\sigma_2^{-1})^m}.
\]
The strand permutation is a power of a three-cycle, so $W_m$ is a knot
exactly when $\gcd(m,3)=1$. Its displayed diagram has $2m$ crossings and
writhe zero. A familiar three-braid identity expresses its Jones polynomial
using reduced Burau matrices \cite{birman1985,alsukaiti2024}:
\begin{equation}
 V_{W_m}(t)=t+t^{-1}+\tr(M(t)^m),\qquad
 M(t)=
 \begin{pmatrix}
  1-t&-t^{-1}\\
  1&-t^{-1}
 \end{pmatrix}.
 \label{eq:weaving-jones}
\end{equation}
Direct calculation gives
\[
 \det M(t)=1,\qquad \tr M(t)=1-t-t^{-1}.
\]
Thus $T_m=\tr(M(t)^m)$ obeys
\[
 T_0=2,\quad T_1=1-t-t^{-1},\quad
 T_{m+1}=(1-t-t^{-1})T_m-T_{m-1}.
\]
Only this elementary recurrence is needed for the specialization audit.

\begin{theorem}[An exact five-color blind family]\label{thm:blind}
At either root of $t^2-3t+1=0$,
\[
 V_{W_m}(t)=3+2(-1)^m.
\]
In particular every odd $m\ge5$ with $3\nmid m$ gives a nontrivial knot
whose exact five-color Jones value equals the unknot value.
\end{theorem}
\begin{proof}
At this specialization, $t+t^{-1}=3$ and $\tr M=-2$.
The trace recurrence has solution $T_m=2(-1)^m$.
Equation~\eqref{eq:weaving-jones} gives the displayed value.
Nontriviality follows either from Theorem~\ref{thm:six} below or from
the determinant formula
\[
 \det(W_m)=L_{2m}-2
   =\left(\frac{3+\sqrt5}{2}\right)^m
    +\left(\frac{3-\sqrt5}{2}\right)^m-2>1
\]
for $m\ge2$ \cite{alsukaiti2024}. Here $L_j$ denotes the Lucas sequence.
For example $W_5$ has determinant $121$.
\end{proof}

This is a collision in a characteristic-zero number field. Both embeddings
have the same value, and every finite-field image in which the specialization
is defined inherits the equality. More precision, a conjugate root, or more
primes at the same $q=5$ point cannot solve it.
The statement is a direct consequence of a classical trace identity; we
make no claim that this particular consequence has not appeared elsewhere.
Its role here is to disprove an overly optimistic completeness inference and
to supply an infinite regression family.

\subsection{Why six colors changes this family}

\begin{theorem}[Every exact $q\ge6$ detects the nontrivial weaving members]
\label{thm:six}
Let $q\ge6$ be an integer and let $t+t^{-1}=q-2$.
For every $m\ge2$ such that $W_m$ is a knot,
$V_{W_m}(t)\ne1$.
\end{theorem}
\begin{proof}
Put $a=q-3\ge3$ and define
\[
 U_0=2,\qquad U_1=a,\qquad U_{m+1}=aU_m-U_{m-1}.
\]
The trace recurrence gives $T_m=(-1)^mU_m$, hence
\begin{equation}
 V_{W_m}(t)=a+1+(-1)^mU_m.
 \label{eq:qweaving}
\end{equation}
We have $U_2=a^2-2>a$. If $U_m>U_{m-1}>0$, then
$U_{m+1}>(a-1)U_m\ge2U_m$. Thus $U_m$ is positive and strictly
increasing for $m\ge1$.
For even $m\ge2$, \eqref{eq:qweaving} is greater than $1$.
For odd $m\ge3$, $U_m>a$ and the expression is less than $1$.
The remaining member $W_1$ has value $1$ and is the unknot.
\end{proof}

At $q=6$, the recurrence starts $2,3,7,18,47,123,\ldots$.
Thus $V_{W_5}=4-123=-119$ and $V_{W_7}=4-843=-839$.
These integer values in a quadratic specialization provide especially simple
independent checks of normalization and color count.

\begin{table}[htbp]
\centering
\begin{tabular}{@{}rrrrr@{}}
\toprule
$m$ & Crossings & Knot? & Exact $q=5$ value & Exact $q=6$ value\\
\midrule
1&2&yes&1&1\\
2&4&yes&5&11\\
4&8&yes&5&51\\
5&10&yes&1&$-119$\\
7&14&yes&1&$-839$\\
8&16&yes&5&2211\\
\bottomrule
\end{tabular}
\caption{Weaving specialization values, computed from the displayed trace
recurrence and independently checked by the transfer routines.}
\label{tab:weaving}
\end{table}

The equality-pattern spaces for $q=5$ and $q=6$ are identical for $f\le5$.
At $f=6$ their sizes are $202$ and $203$, respectively. On small frontiers,
changing the specialization can therefore improve detection with little
change in represented state count. Its worst-case exponential base does
increase, so this observation is not a universal runtime comparison.

\subsection{What this result does and does not decide}

The exact backend defaults to $q=6$ because the $q=5$ family is a proved
structural weakness. A fixed $q=6$ equality is still inconclusive.
Proving that $V_K(2+\sqrt3)=1$ forces the unknot for all knots would imply
the usual Jones unknot-detection conjecture; no such theorem is supplied or
assumed.

Moreover, the baseline's short-braid stage already decides these weaving
examples from their braid representation. Their value here is an adversarial
test of a \emph{filter}, not a newly solved class for end-to-end recognition.
In isolated filter benchmarks that earlier stage is deliberately bypassed
and the bypass is reported. Normal-pipeline benchmarks retain it.
